% Nouveau EGA V (cote 25): archive cote 25, site page 36 (PDF page 37).
% typescript. Draft compared: none.
% AI-transcribed by Claude Opus 5.5 (claude-opus-5-5, adaptive thinking, effort high), Message Batches API, batch ega25b (msgbatch_01A7n5YuXubUPNK1zsepi8RG), 2026-11-13; not yet checked by eye.
\scanpage{25/36}
\mspage{32}

de Z$_n$ , avec codim (Z$_{ni}$,T) = n+k+\unsure{1} et dim f(Z$_{ni}$) = 0 , il existe
une compos ante irréductible T$_j$ de T contenant Z$_{ni}$ telle que
codim (Z$_{ni}$,T$_j$) = n+k+1 et dim f(\xout{Z}T$_j$) = 0.

Lorsque f est quasi-fini donc pour toute partie fermée R de X,
on a dim f(R) = dim R, alors ce critère prend la forme: \inserted[typed]{il n'existe} \xout{pourxtoutxx}
\inserted[typed]{pas de point isolé \xout{dim} z d'un des Z$_n$}
\xout{pointxisoléxxxzxxxdexZx} \xout{Zxxxxxxx} tel que dim$_z$T ( = dim F$_z$ ) soit égal à
n+k+1 \xout{(xàxjoindrexàx8.6.xx)x} Lorsque F est équidimensionnel, \xout{cela}
de dimension d, \xout{celaxsignifiexaussixsimplementxquexxZ} cette condition
est vide si d $\hand{\leqslant}$ k (\xout{et} en effet, nous le savions, car dans c e cas l'hyp
(S$_k$) sur F n'est autre que l'hyp CM), et si d $\hand{\geqslant}$ k+1, elle signifie
que l'ensemble Z$_{d - (k+1)}$ des points de T où la co\unsure{prof}ondeur de F est
$\hand{>}$ d-(k+1) (lequel ensemble est de codimension $\hand{\geqslant}$ d \inserted{i.e. fini} en vertu de
l'hypothèse (S$_k$):sur F) est \uline{vide} , i.e. que \xout{l'ensemblexdes}
l'on a \xout{coprofxFxxxxd-(k+1)}
\[ \text{coprof F } \hand{\leqslant} \text{ d- (k+1)} \]
i.e. vraie prof F $\hand{\geqslant}$ k+1
(alors qu'à priori on a seulement vraie prof F $\hand{\geqslant}$ k comme conséquence
de la propriété (S$_k$) et k $\hand{\leqslant}$ d). \inserted[typed]{Lorsqu'on ne suppose plus F} \deleted{Ce cas particulier mérite d'être}
\deleted{également explicité en corollaire.} équidimensionnel, il reste qu'on
peut exprimer la condition voulue de la façon simple suivante:

\margin{\hand{(8.6.)}} Pour tout \inserted{point fermé} x $\hand{\in}$supp F tel que dim F$_x$ $\hand{\geqslant}$ k+\hand{1}, on a prof F$_x$ $\hand{\geqslant}$ k+1 .

La suffisance se voit aussitôt en faisant z=x, \over{La}{la} nécessité \inserted{se voit} \deleted{aussi,} en
notant que pour tout $\hand{\zeta}$ tel que $\hand{\phi_\zeta}$ soit X$_{k(\hand{\zeta})}$\hand{-}régulier et x $\hand{\in}$ Y$_{\hand{\zeta}}$ ,
on aura dim G$_{\hand{\zeta} x}$ = dim F$_x$ - 1 , prof G$_{\hand{\zeta} x}$ = prof F$_x$ - 1 , donc si
x met en d éfaut la condition plus haut, on aura prof \over{F}{G}$_{\hand{\zeta} x}$ $\hand{<}$k, mais
dim G$_{\hand{\zeta} x}$$\hand{\geqslant}$k, ce qui montre que G$_{\hand{\zeta}}$ ne satisfait pas à (S$_k$) en x; or
l'ensemble d es $\hand{\zeta}$ tels que x$\hand{\in}$ Y$_{\hand{\zeta}}$ est \uline{de codimension} 1 (NB j'ai impli
\margin{|} citement supposé dans le raisonnement que k est alg clos, cas auquel on
se ramène aussitôt). Le critère général précédent doit être \inserted{\unsure{énoncé en cor.}} \deleted{joint à} 8.6
% NOTE: the right edge of the page is cut; the last characters after "8.6" (and the end of "l'hyp", "impli") may be lost
\mspage{36}
